\section{Resolución de tareas}
Hemos resuelto todas las tareas salvo las tres GGH con reducción módulo $n$ (3.01, 3.02 y 3.05),
porque no llegamos a una interpretación clara de qué problema plantean. En total, 12 de nivel~1 (7 GGH y 5
LWE), 10 de nivel~2 (5 GGH y 5 LWE) y 10 de nivel~3 (2 GGH, 5 LWE y 3 Ring-LWE), además de las 10
de Tier~0. Llamamos $k.j$ a la tarea~$j$ del tier~$k$ de GGH, y L$k.j$ y R$k.j$ a las de LWE y
Ring-LWE. Todas las soluciones están comprobadas: en
GGH, que el punto es de la celosía (y en SVP distinto de cero); en LWE, que el secreto deja un error
pequeño en todas las ecuaciones; y en las diez de Tier~0, que el resultado coincide con el del
fichero. Hay además 269 pruebas automáticas que comparan métodos entre sí.

\subsection{GGH}
La Tabla~\ref{tab:ggh} recoge las 14 tareas. En todas el punto de partida (Babai, o el vector más
corto tras LLL) ya era el óptimo salvo en 2.03, donde la enumeración lo mejora, y la enumeración
visita como mucho $n+2$ nodos en dimensión $n$. Las bases públicas
tienen un cociente de Hadamard entre 0{,}007 y 0{,}24, salvo 1.01 y 1.03, que ya eran casi
perpendiculares (0{,}91 y 0{,}99); tras LLL todas superan 0{,}95.

\begin{table}[htb]
\centering
\setlength{\tabcolsep}{5pt}
\begin{tabular}{@{}lcrrl@{}}
\toprule
Tarea & Tipo & $d^2$ & ms & Punto obtenido\\
\midrule
1.01 & CVP & 1 & 0{,}19 & $(36,28,30)$ \\
1.02 & SVP & 17 & 0{,}20 & $(-4,0,-1)$ \\
1.03 & CVP & 16 & 1{,}19 & $(414,152,116,537,502,185)$ \\
1.04 & SVP & 237 & 1{,}35 & $(0,2,2,15,2)$ \\
1.05 & SVP & 226 & 0{,}08 & $(-15,1)$ \\
1.06 & SVP & 229 & 0{,}11 & $(-15,2)$ \\
1.07 & CVP & 5 & 0{,}28 & $(38,-51,84)$ \\
2.01 & SVP & 230 & 0{,}47 & $(0,1,15,2)$ \\
2.02 & SVP & 263 & 1{,}08 & $(-1,2,-1,-16,1)$ \\
2.03 & CVP & 155 & 1{,}26 & $(807,997,549,1528,155)$ \\
2.04 & CVP & 12 & 2{,}16 & $(3842,1730,873,1461,833,3106,1477)$ \\
2.05 & SVP & 263 & 1{,}67 & $(0,-1,0,2,-1,-1,16)$ \\
3.03 & CVP & 8 & 21{,}26 & $(-75,35,2,-76,116,166,52,-233,-154,28,-12,101,119,36,102)$ \\
3.04 & SVP & 250 & 16{,}29 & $(1,0,-15,1,1,-1,0,0,-2,2,0,-2,-1,2,-2)$ \\
\bottomrule
\end{tabular}

\caption{GGH con LLL y enumeración. $d^2$: distancia$^2$ óptima ($\|v\|^2$ en SVP, $\|v-t\|^2$ en
CVP); ms: tiempo de LLL más enumeración, mediana de \Repeticiones{} repeticiones. La dimensión es
el número de coordenadas del punto. En SVP también vale el opuesto del punto.}
\label{tab:ggh}
\end{table}
